% Independent analytical audit, 22 October 2026.
% Incorporatable fragment; uses the notation and TF lemmas of tensor_audit.tex.
\section{The bounded positive-$d$ source union: independent audit}
\label{sec:positive-scale-independent-audit}

\paragraph{Scope and result.}
This audit concerns components 173--176, 193, 196, 198, 200--202 and
205 of the recovered October source index.  Their mathematical content
has been read, including source estimates, all scalar and first-normal
identities, higher-normal restoration, and projective coverage.  No
failure was found in this union at its stated interface: a physical
canonical C1 or C2 source with an actual finite noncommuting contact,
the original native signs and gaps, and coefficients bounded after
normalizing $d=1$.  Entry of an arbitrary supporting functional into
this interface remains a separate theorem.  The conclusions below do
not infer that global entry from a local chart.

The independent script \texttt{positive\_scale\_independent.py} builds
the full parent afresh, using no archive algebra engine.  It checks
34 targeted exact identities, including complete graph scalars,
exceptional first-normal coefficients, the finite-ratio full tail,
all four difficult WN25 conductors, the retained WN37 mass, the three
whole reciprocal core identities, and the RC4 soft pivots derived
directly from contact derivatives.  These checks corroborate algebra.
The following arguments establish the noncomputational points that
such identities cannot decide.

\subsection{A short precise replacement for the critical-center arguments}

All uses of the implicit-function theorem take place in the source
field.  Let $E(c,z)=0$ be a finite polynomial system with real base
point $(c_0,z_0)$ and invertible $z$-Jacobian.  The real local
implicit-function theorem gives open semialgebraic neighborhoods and
a unique semialgebraic solution branch.  Existence and uniqueness in
these fixed neighborhoods are first-order real statements and hence
transfer to every real closed extension.  If $c-c_0$ is infinitesimal,
the resulting source point $z-z_0$ is infinitesimal: for each positive
real neighborhood of $z_0$, shrink the real parameter neighborhood
using continuity and then transfer.  Thus these centers belong to
the original source field and have bounded coordinates; no completion,
convergent power series in the tensor ring, or testing coordinate is
used to construct them.

The scalar sign split in both recovered critical-center terminals
has the following common proof.  Here the equations defining the
branch imply $K_q=0$ on that branch, including the divided $x$
equation when $K_x=U E_x$.  In particular $K_q=0$ at the real
base point.  On the source critical branch put
$k(c,U)=K(q_*(c,U),U)$, with $k(c_0,0)=0$ and
$k_U(c_0,0)=\gamma>0$.  If $\kappa=k(c,0)\geq0$, retain the
nonnegative source term $\kappa U$ and use the center at $U=0$.
If $\kappa<0$, solve
\[
 E(q,U)=0,\qquad K(q,U)+UK_U(q,U)=0.                 \tag{PA1}
\]
At the real base point the augmented Jacobian is block upper
triangular: its first block is the original invertible Jacobian,
its lower-left block is zero, and its last pivot is $2\gamma$.
On the critical branch the second equation is
$g(U)=k(U)+Uk'(U)=0$.  Its derivative is positive in a fixed real
neighborhood and $g(0)=\kappa<0$, so its nearby root has $U_c>0$.
This real local monotonicity assertion transfers.  Evaluate the full
physical parent at the actual source realization with small coordinate
$\sqrt{U_c}$.  Its remaining scalar $\mu$ is nonnegative.  Finally,
ordinary finite polynomial division, at this source center, writes the
fixed-center scalar as $\mu+(U-U_c)^2h$ with $\bar h=\gamma>0$.
This division uses a monic divisor.  It introduces no $U_c^{-1}$ and
no inverse of the independent testing coordinate.

\paragraph{Finite edge (173).}
Normalize the testing event $vW-X+uvY-uZ$ to one; its residue is
$v_0\ne0$, so its inverse is a bounded pure-testing unit.  With
$X=u^2x,Y=ua,Z=ub$, the complete source-center expression is
$4m+u^2K(x,a,b,v,u^2)$ and $K_x=U E_x$.
The divided $x$ pivot at $U=0$ is
$4(1+2n+2\beta)(1+v^2)>0$.
The residue $3$-state complement $M_0(v_0)$ is positive definite.
Its minimized angular profile is a positive real multiple of
$E(v^2)$, where
$E(V)=q_{12}(V-v_0^2)^2$ and $q_{12},v_0^2>0$.
Consequently the full $(a,b,v)$ Hessian is positive definite:
completing the positive $(a,b)$ block leaves the strictly positive
second derivative of this scalar profile.  This is a full Hessian
argument, not an inference from one determinant sign.

The equations $(E_x,K_a,K_b,K_v)=0$ have a block upper triangular
Jacobian with the divided $x$ pivot and that full Hessian.  Their
optimized next coefficient is
\[
 \gamma=\frac{16(1+v_0^2)^2 B(v_0^2)}
                 {v_0^2\det M_0(v_0)}>0.             \tag{PA2}
\]
To see this without a tensor inverse, use the real state Schur
complement.  Event normalization changes the state determinant by
$v^{-2}$.  The scaling of $(X,Y,Z)$ changes its determinant by
$u^8$, while the $K$ state block has determinant
$U\det M_0+O(U^2)$.  The determinant factor
$UE(v^2)+U^2B(v^2)$ therefore gives the minimized $K$ profile.
At $E(v_0^2)=0$ its next coefficient is exactly (PA2); derivatives
of a positive denominator do not contribute there.

Use the actual tensor rows
\[
 X-u^2x_c,\quad Y-ua_c,\quad Z-ub_c,
 \quad u(v-v_c),\quad u^2-U_c.                       \tag{PA3}
\]
The exact state-quadratic expansion has full first coefficients
$u^2E_x,uK_a,uK_b$.  Dividing their differences at the source
center yields only the products printed in (15)--(16) of component
173.  In particular every coupling to the final row has an extra
$u$ or $u^2$, hence infinitesimal coefficient.  The possibly nonzero
couplings among the middle three rows are precisely the positive
$(a,b,v)$ Hessian already proved.  Thus the full residue Gram is
positive definite.  Divided differences of powers of $v^{-1}$ are
finite expressions in the bounded pure units $v^{-1},v_c^{-1}$.
There is no independent-rate assumption in this argument.  Literal
$u=0$ remains an identity in the undivided rows.

\paragraph{Finite-ratio corner (174).}
Normalize $vY-Z=-1$, so $Z=1+vY$.  The opened alpha packet is
exactly $4\alpha(1+u^2)$, including its $u^2$ term.  The formal
center substitution $(X,Y,u,W)=(vx,vy,vt,w)$ gives
$F=4\alpha+v^2K(x,y,t,w,v^2)$.
The positive $(W,x,y)$ complement and
\[
 q_{20}t^4+q_{11}t^2+q_{02}
       =q_{20}(t^2-t_0^2)^2,\qquad q_{20},t_0^2>0
\]
give a strictly positive full $(x,y,t,w)$ Hessian by the same
real Schur argument.  The final pivot is
$16(q_{21}t_0^4+q_{12}t_0^2)/\det M_J(t_0)>0$.
The common construction (PA1) applies.

For the full tail use
\[
 N_1=X-vx_c,\ N_2=Y-vy_c,\ N_3=u-vt_c,
 \ N_4=v(W-w_c),\quad R=v^2-V_c.                   \tag{PA4}
\]
No ratio $u/v$ is needed.  Put $z=W-w_c$ while expanding.  The
full linear part is $v\nabla K\cdot(N_1,N_2,N_3,N_4)$.
Every remaining monomial either has two of $N_1,N_2,N_3$, or has
one of them together with $vz$, or contains $v^2z^2$.
This can also be seen directly from the original event rows:
after normalization the $m$ row is $vW-X-u$, and all other
$W$ appearances are through $uW$; writing $u=vt_c+N_3$ gives
exactly these three possibilities.  Factors $z$ left in
coefficients are infinitesimal.  The independent fresh expansion
checks every monomial of this full tail.  Its residue Gram is
one half of the positive four-variable Hessian.  Couplings to $R$
carry a further $v$.  Hence (PA4) covers every independent rate,
including $v=0$ and unbounded $u/v$.

\subsection{The regular and finite interfaces are sufficient as stated}

For the strong source of 175, $C,A,D,\alpha$ are units.  If $m$
is weak, the real Laurent determinant has strictly positive
$U\partial_U L$ at every interior angle, excluding an interior
zero of a nonnegative function.  If $m$ is a unit, the positive
$(-1,0),(0,-1)$ monomials give a positive angular Hessian in
logarithmic coordinates.  The real $Z=0$ state block is positive,
so the complete five-variable Hessian at a zero is positive.
The source-only implicit center and finite Taylor expansion give
$F=\mu+N^THN$, $\mu\geq0$, $\bar H\succ0$.  This proves the
regular input in 196 independently of a historical A6 label.

The only non-strict finite edge is the $W$ ray at $u=0,v_0\ne0$
with weak $m$.  On this face $q_{10},q_{12},q_{21}>0$ and the
displayed $M_0$ is positive, so (PA2)--(PA3) apply.  The four
literal double corners have the exact positive residue forms
listed in AU9.  Each is manifestly strict on its stated rows:
its positive squares successively force every row to zero.
The complete remainder belongs to the square of those row ideals.
Thus those arguments use neither density nor continuity of SOS
certificates.

For 176, $n,m,C,A,D$ are units and alpha is weak.  The Laurent
radial derivative is positive, excluding an interior zero.
At the sole possible new corner $u=v=0$, the $X,Y$ block is
strict.  In the $Z$ chart the angular Schur complement is PSD;
a nonzero kernel cannot have either angle component zero because
the unit $n,m,A,D$ squares would then force all components zero.
Hence a singular angular block has a one-dimensional kernel at
a finite nonzero ratio.  The positive $W,X,Y$ jet complement and
the positive endpoint and next determinant coefficients verify
every hypothesis of (PA4).  The selected real ratio locates a
source-side jet center, not the independent testing ratio.

\subsection{The weak-alpha leaves: estimates and the exceptional completion}

Components 200--201 first prove $C,A,D$ are units by explicit
forbidden residue faces.  Physicality passes to these real residues
because evaluation at every fixed real tuple and the nonnegative
inequality pass to standard part.  In the determinant balances,
all Laurent coefficients outside the center are nonnegative.
Their bounded sum at contact bounds each term separately.
Writing $U=u_c^2,V=v_c^2$ gives
\[
 U,V=o(1),\qquad m\asymp(U/V)^2(n+V),\qquad
 \alpha\asymp V(U+V).                               \tag{PA5}
\]
For example multiply the radial balance by $V$: its denominator
is $a_3+ma_1V/U$, a positive bounded unit.  This proves the last
comparison, not merely one of its inequalities.  If either $U$
or $V$ had positive residue, the two angular balances force the
other to have positive residue, contradicting weak alpha.
When $n$ is weak, $mV/U,nU/V$ are infinitesimal: split into
bounded or unbounded $U/V$ and use bounded $m$ in the latter case.
These are all source comparisons.

Eliminating the source $X,Y$ block, which has a positive unit
determinant, proves the bounded shape variables before any residue
equation involving them is used.  The resulting critical residue
relations are WA9--WA10.  They imply every displayed hard-block
strictness and the positive pivots
$\mathcal C_0=r_0(1+y_0+4m_0y_0)/2$ and
$\mathcal F_0=r_0^2(1+y_0)/2$ in 201.

For bounded actual $Q=Z_c^2$, the source estimates reduce to
$m\asymp Q^2(n+P)$, $\alpha\asymp P^2$ and
$E\asymp n+P^2$.  The $R,S$ scalar is retained in full;
WN5 and WN6 are exact positive circuit identities.  A fixed real
pivot on $PS$ can be added using WN8 at infinitesimal cost to
the existing positive pivots.  This controls the bounded, rather
than infinitesimal, first-normal coefficients on that row.
The hard residue is WN12, including its nonzero
$(N_X+\bar zN_u)^2$ term.  All terms of hard-normal degree at
least three have infinitesimal coefficient after allocating two
hard rows.  The complementary state chart has the same property.

For unbounded $Q$, one has
$PQ=o(1)$, $m\asymp Q^2(n+P)$ and $\alpha\asymp P^2Q$.
In the $W$ chart a unit $m$ gives the strict ordinary Hessian
WN19.  A weak $m$ gives strict soft rows
$\sqrt m\,v,v^2z,\sqrt\alpha\,z$; the leading first-normal
coefficients divided by $\sqrt m$ are infinitesimal, using
$PQ/\sqrt m=O(\sqrt P)$.  The remaining monomials carry an
extra testing $v$ on the second row.

The complementary unbounded-$Q$ chart is the delicate part.
Put $\epsilon=1/Q$, $\Lambda=PQ$, $\tau=PQ^2$,
$\nu=n/\epsilon^2$, $N=E/\epsilon^2$ as in WN21.
The exact source equations give
\[
 N,\nu,\tau=O(m),\quad \Lambda=O(m\epsilon),\quad
 N=m^2/H_X+\nu H_0r^2/H_X+O(m\Lambda).               \tag{PA6}
\]
No testing parameter occurs here.  If $\nu/m$ has positive
residue, the full $(N_X,\sqrt m A_1)$ block has determinant
$(\nu/m)H_0r^2+O(\Lambda)>0$ in residue.  Otherwise
$\tau/m$ is a positive unit and the completion WN33 is lawful
because it divides only the positive source unit $H_X$.
The signed error $O(m\Lambda)A_1^2$ is infinitesimal on the
six cubic rows by WN25 and the exact identity
$m\Lambda A_1^2=(m/\tau)(\tau\sqrt\epsilon A_1)^2$.

Three details prevent this completion from omitting part of the
parent.  First, the entire additional scalar
$2(m/H_X)A_1P_X$ is allocated using WN26 and WN36; the
$K_1^2$ contribution is infinitesimal either directly, for weak
$m$, or through WN27 when $m$ and hence $\tau$ is a unit.
Second, the exact nonnegative mass
\[
 \{(n+\beta+\tfrac12)v^2+nu^2+
               (n+D+\tfrac12)u^2v^2\}N_X^2          \tag{PA7}
\]
is retained whole; it remains a tensor SOS even after the new
normal is introduced, since each coefficient is a sum of
nonnegative pure-source factors times testing squares.
Third, the full mixed $X$ derivatives are WN38.  Their shift
terms use the third WN25 conductor; every other shifted monomial
has at least two remaining hard normals and infinitesimal
coefficient $A_1$.  Thus the complete higher tail, not just the
scalar restriction, is controlled.  Literal $n=0$ removes every
$\sqrt\nu A_1$ term identically and requires no division by it.

\subsection{The reciprocal leaves and their source pivots}

For C2 (205), the two positive angular balances give
$m\asymp U^2$, $D\asymp V^2$, with $U,V=o(1)$.
Source kernel elimination proves the bounded positive shapes
$S,y,S-y,1-y,R$ of RC2 before they are used as divisors.
One can derive its two soft pivots directly, without importing a
native-coefficient limit.  Substitute into the whole reciprocal
parent
\[
 u=k,\ w=p,\ W=Sp/k,\ X=-Rkp^2,\ Y=yp,\ Z=1,
 \quad m=\mu k^4,\quad D=\delta p^4.                \tag{PA8}
\]
Let $G(k,p)$ denote this graph, taking derivatives before the
last two source substitutions.  State stationarity makes the
graph derivatives agree with the actual angle derivatives.
The normalized source equations $G_k/(2kp^2)=0$ and
$(pG_p/2-G)/p^4=0$ have bounded coefficients after (PA8).
Using the four normalized state equations in residue gives
\[
 C_0=\frac1{2(S-y)},\quad A_0=\frac{1+S-y}{2y},\quad
 \alpha_0=\frac S{2(1-y)},\quad
 n_0=\frac{1+S-y-R}{2R},
\]
and the two angle equations reduce exactly to
\[
 -\mu_0 S^2+\frac{1+S-y}{2}=0,\qquad
 -\delta_0+\frac{S(1+S-y)}2=0.                      \tag{PA9}
\]
Thus both soft pivots are positive units.  The independent
symbolic check verifies these normalized full equations, including
the cancellation of apparently constant terms before division.

The unrestricted tensor rows are RC7.  The source contact makes
the two actual first-normal identities RC11--RC12 hold, with
only the $-C_0N_y\rho$ residue coupling.  In RC14--RC15,
division by $(b-V)^2$ and by $b(b-V)$ is monic polynomial
division.  The coefficients remain bounded because $V$ is
bounded.  The remainders vanish by the source contact jets;
showing that zero never inserts $1/V$ in a coefficient.  Hence
$J_0(0)=D/V^2$ is a conclusion about the already bounded
quotient.  The complete five-row residue is RC17, with its
$N_y,\rho$ determinant $C_0(A_0+\alpha_0+1/2)>0$.
The exact outer core formulas RC20, RC22, RC24 were rebuilt
and independently checked.  Their row Grams have the printed
strict residues and hold at literal projective endpoints.

For C1 (193), the bounds before the graph construction are
substantive and valid.  In the notation CR1--CR8,
$D/P^2$, $\beta U/P^2$ are bounded by the positive angular
balance, and $m/U^2\asymp1+\beta/P$.  This yields
$m/U,D/P,\beta U/P=o(1)$.  The hard source block then gives
unit $r$, bounded $x,y$, and the critical CR7 residues.  Its
full $X$ equation has a numerator consisting of $r\beta$ plus
a positive unit times $P$, modulo $o(P)$; therefore
\[
 -x\asymp P+\beta,\quad r+x\text{ is a positive unit},
 \quad E\asymp\beta+P^2,\quad xQ=o(1).              \tag{PA10}
\]
The scalar balances imply bounded $\delta=D/P^2$,
$\theta=EQ/P$ with $\delta+\theta=\mathcal C+o(1)$;
at least one is a unit.  The full graph scalar and its special
first-normal coefficients were independently rebuilt.

For bounded $Q$, weak beta gives the positive $K$ pivot
$\mathcal F$; unit beta forces $Q$ weak and is used in that
graph only when the testing $z$ is weak.  In the latter case,
CR23--CR25 either add a positive $K$ pivot at infinitesimal
cost or make $K$ itself an infinitesimal combination of circuit
rows.  This correctly handles signed $\mathcal F$.
For unbounded $Q$, $\beta/P=o(1)$, $E/P=o(1)$ and the
exact CR27 identity gives the strict two-row scalar Gram.
The complementary state formulas CR31--CR36 have bounded
coefficients; CR35 controls $aR'$ with infinitesimal coefficients
on the same cubic rows.  Every full first normal is thus
controlled.  The positive $J$ packets are retained whole,
including their possibly unbounded rows.  The hard residues
CR37--CR39 cover the remaining unit-beta states.  All higher
normal terms have bounded coefficients and degree at least three.
No source/test comparison appears in these cases.

For the C1 projective pitchfork (202), the angular balance
forces $D=V^2\delta$, $\beta=V^2\eta$, with bounded
nonnegative $\delta,\eta$ and unit $\delta+\eta$.
The source $W,Y,X$ block PP4 proves bounded centers.
Equations PP9--PP14 are monic divided differences with these
bounded centers.  In particular the division defining $J$
is by the polynomial variable $b$, after its constant term
has been proved zero, not by a testing value.  The entire
parent lies in the five rows PP7.  Their last pivot is
$\bar\delta+\bar\eta\bar c^2>0$.  The other four-row
residue has hard block PP4 and angular Schur pivot
$(1+c^2)^2q_{21}c^2/(4\det H_c)>0$.  This proves full
strictness including unrestricted states.  The projective
complement uses the same complete reciprocal core identities.

\subsection{The source partition and the three-event wall}

The unit/weak source table AU15--AU18 is exhaustive on the
stated canonical interface.  Its forbidden faces follow from
strictly negative central determinant coefficients after all
positive Laurent terms are sent to zero along the specified
real paths.  The transformations AU12--AU13 are full-parent
identities with positive pure-testing norm factors.  For C2,
the remaining balance gives $DV^2\asymp1$, $U^2\asymp m$.
For C1, boundedness of $\beta UV$ ensures that a unit $m$
has unit $U$, while weak $m$ has weak $U$.  In the latter
case bounded $V$ is equivalent to unit $D$.  These implications
give exactly the strong, reflected strong, projective-pitchfork,
and reciprocal leaves, with no rate interval omitted.

Component 198 is independently sound.  Its selected source
cone is the explicit spectrahedral cone
$c\geq0$, $B(c)\succeq0$, because the real Schur complement
gives $\det B=8\Delta$ and
$\Delta=2A(4CD-2Dd-d^2)-d^2(d+2C+2D)$.
The strictly feasible real point bounds all dual representing
matrices by their trace; hence the displayed dual cone is
closed over the reals and finite-dimensional separation gives
CAP10.  The resulting finite first-order statement transfers
to the testing field.  Separate PSD factorizations then give
the tensor certificate; no non-Archimedean compactness argument
is used.  At any finite noncommuting contact the selected
parent would be positive definite if $\Delta\geq0$, so the
full parent cannot vanish there.  All eight wall inequalities
are therefore strictly negative at a finite contact, by the
exact cyclic event symmetry.  This excludes the advertised
wall intersections and supplies the total-scale dichotomy;
it does not supply arbitrary canonical entry.

\paragraph{Proof-use conclusion.}
The bounded positive-$d$ union can be retained in a refined
proof with one critical-center lemma, one monic-division lemma,
one full-normal allocation rule, and the explicit circuit and
source identities above.  Historical audit receipts and
repeated theorem wrappers are unnecessary.  The native finite
contact hypothesis, its sign/gap interface, and the weak-total-
scale C1/C2 branches must remain explicit separate dependencies.
